% SOURCE_STATEMENT: "Black-hole stability conjecture: do sufficiently small perturbations of sub-extremal Kerr initial data for the vacuum Einstein equations evolve toward a nearby Kerr black hole?"
% Preparation and final self-audit: CONTRIBUTING.md prompts C and D; source/audit artifacts retained in the local physics_statements_2026_10_11 batch.
% !TeX program = lualatex
\documentclass[11pt]{article}
\usepackage[margin=25mm]{geometry}
\usepackage{fontspec}
\setmainfont{Latin Modern Roman}
\usepackage{amsmath,amssymb,mathtools}
\usepackage{unicode-math}
\setmathfont{Latin Modern Math}
\usepackage{xurl,hyperref}
\hypersetup{colorlinks=true,urlcolor=blue}
\setcounter{secnumdepth}{0}
\setlength{\parindent}{0pt}
\setlength{\parskip}{5pt}
\setlength{\emergencystretch}{3em}
\newcommand{\sref}[2]{\hyperlink{source:#1:#2}{[S#2]}}
\newcommand{\cataloguescope}[1]{\paragraph{Recorded target.}#1}
\newcommand{\R}{\mathbb R}
\newcommand{\C}{\mathbb C}
\newcommand{\N}{\mathbb N}
\newcommand{\Z}{\mathbb Z}
\newcommand{\Q}{\mathbb Q}
\newcommand{\F}{\mathbb F}
\newcommand{\eps}{\varepsilon}
\begin{document}
\section{Nonlinear Stability of Subextremal Kerr Black Holes}\label{statement}

\subsection{Definitions and mathematical statement}

\paragraph{Shared notation.}$\mathbb N=\{1,2,\ldots\}$; the standard number fields have their usual meanings. All additional conventions are stated locally.


\paragraph{Definitions and conventions.}
The vacuum Einstein equations are $\operatorname{Ric}(g)=0$ for a Lorentzian spacetime metric $g$. Kerr solutions are the stationary rotating black-hole family with mass $M>0$ and rotation parameter $a$; subextremality means $|a|<M$. Initial data consist of a Riemannian metric and second fundamental form satisfying the vacuum constraint equations. Nonlinear stability concerns the maximal evolution of perturbed data, including its exterior geometry and late-time convergence, with final Kerr parameters allowed to differ from the initial parameters.

\paragraph{Target retained from the canonical source catalogue.}
\paragraph{Catalogue wording retained for traceability.}
Black-hole stability conjecture: do sufficiently small perturbations of sub-extremal Kerr initial data for the vacuum Einstein equations evolve toward a nearby Kerr black hole?

\paragraph{Source scope and qualification.}
The catalogue does not specify the perturbation norm, gauge or convergence topology. Those are parts of the stability formulation in the cited report; the target is the full subextremal family, rather than a silently restricted small-rotation theorem.
\subsection{Short English statement}
Do all sufficiently small perturbations of a nonextremal rotating vacuum black hole settle to another Kerr black hole?
\subsection{Sources}
\begin{itemize}
\item[\hypertarget{source:statement:1}{S1}] ULAM.AI and the UnsolvedMath Contributors, UnsolvedMath, catalogue entry. Catalogue curation under CC BY 4.0; definitions and English exposition added locally.
\url{https://huggingface.co/datasets/ulamai/UnsolvedMath}.
\item[\hypertarget{source:statement:2}{S2}] Nonlinear Evolution Problems (2012), Oberwolfach Reports, DOI 10.4171/owr/2012/26, pp. 1563--1637. Primary report accessed 2026-10-11; relevant PDF pages 11, 12.
\url{https://doi.org/10.4171/owr/2012/26}.
\end{itemize}
\label{end:statement}
\end{document}
